\section{Background and the objective family}
\label{sec:background}

\subsection{Diffusion training and notation}
\label{sec:notation}

A denoising diffusion model \citep{ho2020denoising} with $T$ timesteps and
noise schedule $\beta_1,\dots,\beta_T$ sets
$\bar\alpha_t = \prod_{s \le t}(1-\beta_s)$ and corrupts a datum $x_0$ as
\begin{equation}
  x_t = \sqrt{\bar\alpha_t}\,x_0 + \sigma_t\,\epsilon, \qquad
  \sigma_t = \sqrt{1-\bar\alpha_t} \in (0,1), \qquad
  \epsilon \sim \mathcal{N}(0, I).
  \label{eq:forward}
\end{equation}
We call $\sigma_t$ the noise level ($\sigma_t \to 1$ is the high-noise end of
the schedule) and write
$\mathrm{SNR}_t = \bar\alpha_t/\sigma_t^2 = (1-\sigma_t^2)/\sigma_t^2$ for the
signal-to-noise ratio. A network $\hat\epsilon_\theta(x_t, t)$ is trained to
predict $\epsilon$ under
\begin{equation}
  \mathcal{L}_{\mathrm{simple}}
  = \mathbb{E}_{x_0,\epsilon,t}\big[\operatorname{ms}\big(\hat\epsilon_\theta(x_t,t) - \epsilon\big)\big],
  \qquad
  \operatorname{ms}(z) = \lVert z \rVert_2^2 / \dim z ,
  \label{eq:simple}
\end{equation}
with $t$ uniform on $\{1, \dots, T\}$. The prediction defines a clean-image
estimate $\hat x_0 = (x_t - \sigma_t \hat\epsilon_\theta)/\sqrt{\bar\alpha_t}$.
Substituting Eq.~\eqref{eq:forward} gives, for the residual
$r = \hat\epsilon_\theta(x_t,t) - \epsilon$ of one sample,
\begin{equation}
  r = -\sqrt{\mathrm{SNR}_t}\,(\hat x_0 - x_0),
  \qquad
  \operatorname{ms}(r) = \mathrm{SNR}_t \operatorname{ms}(\hat x_0 - x_0).
  \label{eq:residual}
\end{equation}
Every loss is averaged over channels and over the minibatch, and $N$ is the
number of entries per channel. Two consequences of Eq.~\eqref{eq:residual}
are used throughout. First, the regression target $\epsilon$ is white noise
but the residual is not: it is the model's clean-image error rescaled by
$\sqrt{\mathrm{SNR}_t}$, so it carries image structure and its spectrum
depends on the noise level. For stationary Gaussian data with power spectrum
$P$, the optimal predictor's residual has variance
$\bar\alpha_t P(u)/(\bar\alpha_t P(u) + \sigma_t^2)$ at frequency $u$
(Proposition~\ref{prop:residual}; all propositions are proved in
Appendix~\ref{app:theory}): nearly white
at low noise, and increasingly shaped like the data spectrum, concentrated at
low frequencies for natural images, as the noise level grows. Second, a
timestep weight $\kappa_t$ on the $\epsilon$-loss is a weight
$\kappa_t\,\mathrm{SNR}_t$ on the clean-image error, so
Eq.~\eqref{eq:simple} weights that error by $\mathrm{SNR}_t$.
Min-SNR-$\gamma$ \citep{hang2023efficient} sets
$\kappa_t = \min(\mathrm{SNR}_t, \gamma)/\mathrm{SNR}_t$ ($\gamma = 5$
throughout), lowering the weight of the low-noise timesteps, those with
$\mathrm{SNR}_t > \gamma$; raw SNR weighting, $\kappa_t = \mathrm{SNR}_t$,
raises it without bound.

\subsection{The objective family}
\label{sec:family}
\label{sec:family-def}

A sub-band objective fixes a linear transform $W$ whose output coefficients are
partitioned into bands $b \in \mathcal{B}$, band $b$ holding $n_b$
coefficients, and weights the bands by non-negative functions of the noise
level (Section~\ref{sec:family-ffl} admits weights that also depend on the
residual):
\begin{equation}
  \mathcal{L}_{W,\omega}
  = \mathbb{E}_{x_0,\epsilon,t}\Big[\sum_{b \in \mathcal{B}} \omega_b(\sigma_t)
      \operatorname{ms}\big((W\hat\epsilon_\theta)_b - (W\epsilon)_b\big)\Big]
  = \mathbb{E}_{x_0,\epsilon,t}\Big[\sum_{b \in \mathcal{B}} \omega_b(\sigma_t)
      \operatorname{ms}\big((W r)_b\big)\Big].
  \label{eq:family}
\end{equation}
The one-level 2\nobreakdash-D discrete wavelet transform (DWT) \citep{mallat1989theory}
splits each channel into an approximation band $\mathrm{LL}$ and horizontal,
vertical and diagonal detail bands $\mathrm{LH}$, $\mathrm{HL}$ and
$\mathrm{HH}$ at half resolution; with Haar filters and even image sizes it
is orthonormal and each band holds $N/4$ coefficients. The 2\nobreakdash-D discrete
Fourier transform (DFT) $F$ is applied per channel with orthonormal scaling,
so it is unitary.

By the linearity of $W$ the weighted bands are bands of the residual, and by
Eq.~\eqref{eq:residual} the term inside the expectation equals
\begin{equation}
  \mathrm{SNR}_t \sum_{b \in \mathcal{B}} \omega_b(\sigma_t)
  \operatorname{ms}\big((W(\hat x_0 - x_0))_b\big),
  \label{eq:family-x0}
\end{equation}
a band weighting of the clean-image error multiplied by the timestep factor
$\mathrm{SNR}_t$ that $\mathcal{L}_{\mathrm{simple}}$ already carries. How
that error's energy divides among the bands changes with $\sigma_t$
(Section~\ref{sec:notation}), so a weight that depends on $\sigma_t$ acts on
two things at once: the frequency content of the error and the relative
weight of timesteps. Proposition~\ref{prop:confound} splits the loss exactly
into a residual-independent timestep weight on the pixel loss and a zero-sum
band term, which also reweights timesteps through the residual's band
energies; Section~\ref{sec:confound} examines the consequences for AWWL.

\label{sec:family-parseval}
\label{sec:family-info}
If $W^{*}W = I$, with $W^{*}$ the adjoint, that is, if $W$ is an isometry (as
are the orthonormal DWT with zero or periodic extension and the unitary DFT,
but not a DWT with symmetric extension, biorthogonal filters or unconstrained
learned filters), then for every residual
\begin{equation}
  \sum_{b \in \mathcal{B}} \big\lVert (W r)_b \big\rVert_2^2 = \lVert r \rVert_2^2 ,
  \label{eq:parseval}
\end{equation}
so the weights $\omega_b = n_b/N$ turn Eq.~\eqref{eq:family} into
$\mathcal{L}_{\mathrm{simple}}$ exactly (checked to single precision by
\texttt{scripts/verify\_loss\_math.py}). The condition is $W^{*}W = I$, not
equal band sizes; for the one-level Haar DWT this unweighted limit is
$\omega_b = 1/4$.

If the weights are strictly positive and depend on the timestep only, then
for any injective $W$ the population minimiser of Eq.~\eqref{eq:family} is
the conditional mean $\mathbb{E}[\epsilon \mid x_t, t]$, as for
$\mathcal{L}_{\mathrm{simple}}$ (Proposition~\ref{prop:minimiser}). In
this sense such a weighting adds no information about the target; it acts
only through how a model of finite capacity trades errors between bands and
timesteps and through the gradients that the optimiser receives
(Propositions~\ref{prop:confound} and~\ref{prop:scale}). The statement fails
for weights that depend on the residual (FFL) or on the target (the spatially
adaptive arm of Section~\ref{sec:extensions}; remarks after
Proposition~\ref{prop:minimiser}).
Equation~\eqref{eq:parseval} holds for a fixed residual and does not constrain
which residuals training reaches, so whether a change of emphasis helps is
empirical.

\subsection{The audited objective}
\label{sec:family-awwl}

AWWL, introduced in our conference paper \citep{thao2026awwl}, uses the
one-level Haar DWT (zero extension) with the rational weighting
\begin{equation}
  w_{\mathrm{LL}}(\sigma) = \alpha\,\lambda_p(\sigma), \qquad
  w_{\det}(\sigma) = (1-\alpha)\big(1 - \lambda_p(\sigma)\big), \qquad
  \lambda_p(\sigma) = \frac{\sigma^{p}}{\sigma^{p} + (1-\sigma)^{p}} .
  \label{eq:weights}
\end{equation}
Here $\alpha \in [0,1]$ sets the balance between the approximation band and the
detail bands and $p \ge 0$ how sharply that balance follows $\sigma$. For
$p > 0$, $\lambda_p$ increases from $0$ at $\sigma = 0$ to $1$ at $\sigma = 1$,
so the approximation band dominates at high noise. The published settings are
$\alpha = 0.2$, $p = 1$ on CIFAR-10 and $\alpha = 0.8$, $p = 2$ for DreamBooth
fine-tuning. In the released implementation, which produced the published
results and is audited here, the three detail bands
$\mathcal{D} = \{\mathrm{LH},\mathrm{HL},\mathrm{HH}\}$ are averaged,
whereas the equation printed in the conference paper (its Eq.~7) sums them:
\begin{linenomath}
\begin{align}
  \mathcal{L}_{\mathrm{AWWL}}
  &= \mathbb{E}\Big[w_{\mathrm{LL}}(\sigma_t)\operatorname{ms}(r_{\mathrm{LL}})
    + \frac{w_{\det}(\sigma_t)}{3}
      \sum_{b \in \mathcal{D}} \operatorname{ms}(r_b)\Big],
  \qquad r_b = (W r)_b ,
  \label{eq:audited}\\
  \mathcal{L}_{\Sigma}
  &= \mathbb{E}\Big[w_{\mathrm{LL}}(\sigma_t)\operatorname{ms}(r_{\mathrm{LL}})
    + w_{\det}(\sigma_t)
      \sum_{b \in \mathcal{D}} \operatorname{ms}(r_b)\Big].
  \label{eq:summed}
\end{align}
\end{linenomath}
We call the averaged form, Eq.~\eqref{eq:audited}, AWWL; `published' refers
below only to the conference paper's printed equation, settings and results.
Sums over bands such as $\sum_b \omega_b$ and $\sum_b \omega_b^2$ refer to
the per-band weights of AWWL,
$\omega = (w_{\mathrm{LL}}, w_{\det}/3, w_{\det}/3, w_{\det}/3)$. The
summed-detail variant, Eq.~\eqref{eq:summed}, triples the ratio of detail to
approximation weight at every $\sigma$, so it is a different objective, not a
rescaling of AWWL; Section~\ref{sec:replication} reports it, and arms built on
it are identified where they appear. The total weight of AWWL is
\begin{equation}
  S(\sigma) = \sum_b \omega_b(\sigma) = w_{\mathrm{LL}}(\sigma) + w_{\det}(\sigma)
  = (1-\alpha) + (2\alpha - 1)\,\lambda_p(\sigma) ,
  \label{eq:total}
\end{equation}
which is affine in $\lambda_p$, runs from $1-\alpha$ ($\sigma \to 0$) to
$\alpha$ ($\sigma \to 1$), and is constant only at $\alpha = 1/2$ (for
$p > 0$). The unweighted limit of AWWL is $w_{\mathrm{LL}} = 1/4$,
$w_{\det} = 3/4$; Proposition~\ref{prop:minimiser} covers both forms for
$0 < \alpha < 1$.

\subsection{A second member: Focal Frequency Loss}
\label{sec:family-ffl}

Focal Frequency Loss (FFL) \citep{jiang2021focal} takes $W = F$ with one band
per frequency and weights each frequency by the current magnitude of its own
error:
\begin{equation}
  \mathcal{L}_{\mathrm{FFL}}
  = \mathbb{E}\Big[\frac{1}{N} \sum_{u} w_u \lvert R_u \rvert^2\Big],
  \qquad R = F r,
  \qquad w_u = \operatorname{sg}\bigg[\frac{\lvert R_u \rvert^{a}}{\max_{u'} \lvert R_{u'} \rvert^{a}}\bigg],
  \label{eq:ffl}
\end{equation}
where the maximum runs over all frequencies and channels of the sample,
$\operatorname{sg}$ is the stop-gradient operator, and $a = 1$ as in
\citet{jiang2021focal}; with $a = 0$, Eq.~\eqref{eq:ffl} is the pixel MSE by
Eq.~\eqref{eq:parseval}. FFL differs from AWWL in three respects: its basis is
global rather than spatially localised; its weights depend on the residual
rather than on $\sigma_t$, so Proposition~\ref{prop:minimiser} does not
apply; and it weights timesteps implicitly, since the ratio of the loss to
the pixel MSE,
$\Lambda_{\mathrm{FFL}}(r) = \sum_u w_u \lvert R_u \rvert^2 / \sum_u \lvert R_u \rvert^2$,
is unchanged when $r$ is rescaled but depends on the spectral shape of the
residual, which changes with $\sigma_t$. \citet{jiang2021focal} add FFL to an
existing spatial-domain loss; here it is the sole objective on the
$\epsilon$-residual (Section~\ref{sec:second-member}).
